% !PRINTED 416
\weberpage{416}
%<<CONT
no greater common divisor than 2. If $x$ and $y$ have the greatest
common divisor 2, then $b$ and $c$ are even, and $a$, $b$, $c$,
$\partial$ are then and only then relatively prime, when $a$ and
$\partial$, and hence also $n = a\,\partial - b\,c$, are \gesperrt{odd}.

If we disregard the uninteresting case $n=1$, equation (8) has no
solution satisfying these conditions in which $x=0$. But if we change
the signs of $x$, $y$ at the same time, then by (10) $a$, $b$, $c$, % [sic] German source reference (10); (7) is mathematically expected.
$\partial$ pass over into $-a$, $-b$, $-c$, $-\partial$, and equation
(1) is not altered. The transformation
$\begin{pmatrix} a, & b\\ c, & \partial \end{pmatrix}$ therefore
remains unchanged.

For brevity we now call a solution of equation (8) a
proper solution, and $n = \frac{1}{4}(y^2-D\,x^2)$ a
\gesperrt{proper representation} of $n$, if $x$ and $y$ satisfy the
following conditions:

\gesperrt{$x$ is positive.}

\gesperrt{The greatest common divisor of $x$ and $y$ is 1 or 2.}

\gesperrt{If it is $=2$, then $n$ is odd.}

If now $\omega$ satisfies a given equation (3), then by every system
$\begin{pmatrix} a, & b\\ c, & \partial \end{pmatrix}$ a proper
solution of (8) is uniquely determined. For by (5) $x$ is the greatest
common divisor of $b$, $c$, $a-\partial$, and $y$ is determined by
(6); and conversely, from a proper solution of (8) the system $a$,
$b$, $c$, $\partial$ is completely determined by (7). But the question
is still to be decided whether two \gesperrt{different} proper solutions $x$, $y$
and $x'$, $y'$ can lead to equivalent systems
$\begin{pmatrix} a, & b\\ c, & \partial \end{pmatrix}$,
$\begin{pmatrix} a', & b'\\ c', & \partial'
\end{pmatrix}$. By equivalent systems are here to be understood,
according to \S\,28, \S\,53, two such for which
\weq{9}{\begin{pmatrix} a', & b'\\ c', & \partial' \end{pmatrix} = \begin{pmatrix} \alpha, & \beta\\ \gamma, & \delta \end{pmatrix}\begin{pmatrix} a. & b\\ c, & \partial \end{pmatrix}} % [sic] source prints a full stop after a in the rightmost matrix.
holds, when $\begin{pmatrix} \alpha, & \beta\\ \gamma, & \delta
\end{pmatrix}$ is a \gesperrt{linear} transformation, and hence
\weq{10}{\alpha\,\delta - \beta\,\gamma = 1.}
If we compose both sides of equation (9) on the right with
$\begin{pmatrix} \partial, & -b\\ -c, & a \end{pmatrix}$, there
results:
%CONT>>
